\subsection{AQ-67-NONHOLO: the natural complex-analytic formulation fails}
Give the labelled geodesic-cut $(0,4)$ space its natural complex
structure, namely the covering
\[
 \mathcal C=\mathcal T(S_{0,4})/\langle D_\alpha\rangle
\]
of the moduli space of four labelled punctures, where $\alpha$ is
the chosen separating curve and $D_\alpha$ is its Dehn twist.
Equivalently, this is the component of the cover that remembers
the isotopy class of the cut. The two sides and all punctures
remain labelled. Fenchel--Nielsen coordinates upstairs are
$(\ell,\tau)\in\mathbb R_{>0}\times\mathbb R$, with
$D_\alpha(\ell,\tau)=(\ell,\tau+\ell)$; downstairs the
normalized angle is $\theta=2\pi\tau/\ell$ modulo $2\pi$.
The canonical capping map of AQ-67-EXACT is not holomorphic
from this complex manifold to $\mathbb C^\times$.
Thus its everywhere invertible real tangent map does not make
it a complex-analytic or $\mathbb C$-algebraic \emph{\'etale}
map. This refutes that stronger interpretation of 67/077--082
for the natural Teichm\"uller complex structure already in
the labelled $(0,4)$ sector.

\begin{proof}
Suppose it were holomorphic. By AQ-67-EXACT it is a bijection
with nonzero real Jacobian onto the punctured disk
$\mathcal D=\{0<|q|<16\}$. The holomorphic inverse function
theorem then makes it a biholomorphism with $\mathcal D$.
Conjugating the rotations $q\mapsto e^{is}q$ by it gives
holomorphic automorphisms of $\mathcal C$. Their lifts to
$\mathcal T(S_{0,4})$, continuously chosen to equal the identity
at $s=0$, are the explicit global diffeomorphisms
\[
 H_s(\ell,\tau)=
       \left(\ell,\tau+\frac{s\ell}{2\pi}\right),
                      \qquad s\in\mathbb R.                  \tag{AQ.67.11}
\]
This follows directly from the tensor's unit twist phase;
any choice of zero-angle section changes $\theta$ by a function
of $\ell$ and leaves (AQ.67.11) unchanged. Since the quotient
projection is a holomorphic covering, every $H_s$ is
holomorphic as well.

Let $\omega_{\mathrm{WP}}$ and $g_{\mathrm{WP}}$ be the
Weil--Petersson K\"ahler form and metric. Wolpert (1983),
Lemma 4.1, p.~229, says that the arclength twist vector field
$t_\alpha=\partial/\partial\tau$ has Hamiltonian potential
$-\ell$, in his normalization:
\[
 \iota_{t_\alpha}\omega_{\mathrm{WP}}=-d\ell.
\]
The vector field of (AQ.67.11) is
$V=(\ell/2\pi)t_\alpha$, so
\[
 \iota_V\omega_{\mathrm{WP}}
   =-\frac{\ell}{2\pi}d\ell
   =-d\!\left(\frac{\ell^2}{4\pi}\right).
\]
Cartan's formula and $d\omega_{\mathrm{WP}}=0$ imply
$\mathcal L_V\omega_{\mathrm{WP}}=0$; hence every global
$H_s$ preserves $\omega_{\mathrm{WP}}$. A different fixed
normalization of the Weil--Petersson form only multiplies
these identities by a nonzero constant and gives the same
conclusion. A holomorphic map preserving the K\"ahler form
preserves the metric, since that metric is recovered from
the form and the complex structure. Thus every $H_s$ would
be a Weil--Petersson isometry of $\mathcal T(S_{0,4})$.

Brock and Margalit, \emph{Weil--Petersson isometries via the
pants complex}, Main Theorem, p.~2 of arXiv:math/0412499v1,
prove that every such isometry is induced by an element of
the extended mapping class group; their theorem explicitly
includes $S_{0,4}$. That group is countable. On the other
hand (AQ.67.11) gives distinct maps for every two distinct
real numbers $s$, because $\ell>0$ and the upstairs twist
coordinate is real, not taken modulo $\ell$. This is an
uncountable family of isometries, a contradiction.
\end{proof}
